\textbf{结论名称：} 平面点绕任意中心旋转的坐标公式（V040）

\textbf{英文名称：} rotation\_about\_arbitrary\_center\_coordinates

\textbf{适用条件：} 在平面直角坐标系中，点 $P(x_1,y_1)$ 绕固定点 $B(a,b)$ 旋转角 $\theta$ 到达点 $Q(x_2,y_2)$；\textbf{逆时针取正角，顺时针取负角}，旋转不改变 $BP$ 的长度。

\textbf{变量说明：} 令 $u=x_1-a,\ v=y_1-b$，即从 $B$ 到 $P$ 的水平、竖直变化量；$\theta$ 可取任意角度。

\textbf{核心公式：}
\[
\boxed{\begin{aligned}
x_2&=a+u\cos\theta-v\sin\theta,\\
y_2&=b+u\sin\theta+v\cos\theta.
\end{aligned}}
\]
其中 $\theta<0$ 代表顺时针旋转；直接代入 $-\theta$ 也可。

\textbf{识别信号：} 绕某点旋转、已知角度及旋转中心、求对应点坐标；或已知两组相对坐标、反求旋转角与方向。

\textbf{使用场景：} 三角形旋转、正方形与正三角形坐标、45°与60°特殊角、点轨迹、解析几何消元、垂直等长的90°特例（V039）。

\textbf{简单例子：} $B(2,-1), P(5,3)$，则 $\overrightarrow{BP}=(3,4)$。逆时针转 $90^\circ$ 后，$\overrightarrow{BQ}=(-4,3)$，故 $Q=(-2,2)$。逆时针转 $45^\circ$ 后，$Q=(2-\tfrac{\sqrt2}{2},\,-1+\tfrac{7\sqrt2}{2})$。
